\chapter{Glossary of XNA and CNA Terms}
\label{ch:glossary}

\begin{description}
\item[ASCII] A removed public renderer identity. Its reusable quantization and glyph-atlas logic
  now lives in the renderer-neutral \cnaclass{CNA::Graphics::AsciiPostProcessEffect} extension
  under \texttt{modules/graphics-ext}; the former renderer's architecture is retained as history
  in Chapter~\ref{ch:canvas-ascii-backends}.
\item[BGFX] One of CNA's 46 public renderer identities; a third-party rendering library integrated
  directly via its native API. See Chapter~\ref{ch:bgfx-backend}.
\item[CANVAS] CNA's Emscripten-only, HTML \texttt{<canvas>}-based 2D graphics renderer. See
  Chapter~\ref{ch:canvas-ascii-backends}.
\item[CNA] The project this book is about: a C\texttt{++}23 reimplementation of the XNA~4.0
  programming model, built on SDL3 with a pluggable graphics-renderer layer.
\item[ContentManager] The central XNA class resolving an asset name to real loaded content,
  through a three-way resolution order this book returns to repeatedly: a real, compiled
  \texttt{.xnb} file first, then a \texttt{.cnj} sidecar, then a direct native-format load. See
  Chapter~\ref{ch:content-and-assets}.
\item[DIRECTX9 / D3D9] \texttt{DIRECTX9} is CNA's public selector for its Direct3D~9
  renderer, and Direct3D~9 is the API this book's own zero-tolerance
  oracle methodology (see \textbf{Oracle (D3D9)} below) is built around, since a genuine XNA~4.0
  runtime itself targets Direct3D~9. See Chapter~\ref{ch:direct3d-backends}.
\item[DIRECTX11 / D3D11] \texttt{DIRECTX11} selects CNA's Direct3D~11 renderer, sharing a common \texttt{D3DCommon} foundation
  with D3D12 but implementing mip-chain generation via one native, whole-resource
  \texttt{GenerateMips()} call where D3D12 cannot. See Chapter~\ref{ch:direct3d-backends}.
\item[DIRECTX12 / D3D12] \texttt{DIRECTX12} selects CNA's Direct3D~12 renderer; unlike D3D11 it has no native convenience call
  for mipmap generation and instead implements it as an explicit, face-scoped CPU
  readback/downsample/upload loop. See Chapter~\ref{ch:direct3d-backends}.
\item[Descriptor heap] A Direct3D~12-owned table of resource-view descriptions rather than the
  resources themselves. CNA's D3D12 renderer keeps separate RTV, DSV, CBV/SRV/UAV, and sampler
  heaps, currently using fixed-capacity bump allocation with no slot reclamation during a device
  lifetime; exhaustion is therefore a CNA implementation limit, not a Direct3D~12 resource-memory
  limit. See Chapter~\ref{ch:direct3d-backends}.
\item[FREEDIRECT] CNA's graphics renderer fronting the sibling \texttt{free-direct} project;
  a 2D-only renderer in the same spirit as \texttt{SDL\_RENDERER}. The former \texttt{DX3}
  selector was renamed; current \texttt{DIRECTX3} denotes a different, genuine DirectX~3
  implementation. See Chapter~\ref{ch:dx3-freedirect-backend}.
\item[EasyGL profiles] \texttt{OPENGLES2}, \texttt{OPENGLES3}, \texttt{OPENGL33},
  \texttt{WEBGL1}, and \texttt{WEBGL2} are distinct public identities sharing CNA's internal
  EasyGL factory and the sibling \texttt{easy-gl} project. \texttt{EASYGL} is not an accepted
  selector. See Chapter~\ref{ch:easygl-backend}.
\item[Effect] The XNA base class for a shader program bound to a \cnaclass{GraphicsDevice}.
  CNA ships five stock effects and renderer-specific runtime-compiled custom
  \texttt{ShaderEffect} path --- see the shader/.fx gap chapter, Chapter~\ref{ch:shader-fx-gap}.
\item[ENet] The third-party, vendored UDP networking library backing CNA's real Networking
  namespace implementation, wrapped by an internal \texttt{ENetBackend}. The same library
  \texttt{free-direct}'s own optional DirectPlay transport is built on. See Chapter~\ref{ch:networking}.
\item[\texttt{EXT} suffix] CNA's naming convention for a method or property bolted onto an
  otherwise faithfully-XNA-shaped class, kept visually distinct from the real XNA members
  around it --- for example \cnaclass{GraphicsDevice}'s parameterless, headless-friendly
  constructor, or \texttt{MinimizeEXT()} / \texttt{RestoreEXT()} on \cnaclass{GameWindow}, which
  XNA has no minimize/restore API for at all. A sibling of \texttt{CNAEXT}: both mark non-XNA
  additions, but \texttt{EXT} specifically marks an addition to an otherwise-real XNA class,
  where \texttt{CNAEXT} marks a whole declaration with no XNA counterpart. See Chapter~\ref{ch:game-loop}.
\item[FNA] An open-source, C\#, SDL2-based reimplementation of the XNA~4.0 API, maintained
  independently of CNA. FNA is CNA's primary behavioral reference and the subject of CNA's
  own differential-testing infrastructure --- not the same thing as the real Microsoft
  XNA~4.0 specification, a distinction this book returns to repeatedly (most sharply in the
  Media chapter's \cnaclass{Song} finding, Chapter~\ref{ch:media-system}).
\item[Fence (D3D12)] A GPU-completion counter used to tell when submitted Direct3D~12 command
  work has finished. CNA records one previous fence value for each of its two command-allocator
  slots and waits before reusing that slot, rather than waiting for the newly submitted frame
  itself; its simpler off-screen draw helpers intentionally wait synchronously for deterministic
  readback. See Chapter~\ref{ch:direct3d-backends}.
\item[GamerServices] The XNA namespace covering Xbox-Live-era player profile, achievement,
  leaderboard, and system-UI concepts. CNA implements it with real local persistence where
  FNA does not implement it at all. See Chapter~\ref{ch:gamerservices}.
\item[Golden-image test] CNA's own visual-regression testing pattern: a real rendered frame,
  captured via \texttt{SaveBackBufferScreenshotEXT} / \cnaclass{Texture2D}\texttt{::SaveAsPng}, is
  compared pixel-for-pixel against a checked-in reference (``golden'') PNG, with a
  \texttt{CNA\_UPDATE\_GOLDEN} escape hatch to intentionally regenerate the reference when a
  change is deliberate. This book's own real \texttt{SOFTWARE}-renderer screenshots are built on
  the same infrastructure. See Chapter~\ref{ch:graphicsdevice}.
\item[HEADLESS] A CNA graphics renderer that never rasterizes a real pixel at all --- it only
  tracks and reports the last \texttt{Clear()} color and similar logic-only state. Used for
  fast unit testing and unable to produce a real screenshot for that reason. See
  Chapter~\ref{ch:graphicsdevice}.
\item[GraphicsCapability] The small enum \cnaclass{GraphicsDevice::SupportsCapability} takes,
  letting game code ask whether the active renderer supports a feature (3D rendering,
  occlusion queries, and so on) before attempting it, rather than attempting it and catching
  the resulting exception. Not every value is yet wired to answer honestly on every renderer:
  \cnaclass{MultipleRenderTargets}, \cnaclass{DepthStencilBuffer}, and
  \cnaclass{CustomEffects} have confirmed false positives, and several device-dependent limits
  cannot be represented by one Boolean. See Chapter~\ref{ch:renderer-contract} and
  Appendix~\ref{ch:feature-matrix} for the qualified matrices.
\item[GraphicsDevice] The central class every rendering operation in CNA passes through;
  covered in Chapter~\ref{ch:graphicsdevice}.
\item[.cnj] CNA's own permanent, JSON-based content sidecar format --- ``CNA content JSON,''
  originally and confusingly named \texttt{.cnb}. Tried after \texttt{.xnb} and a direct
  native-format load in \cnaclass{ContentManager}'s resolution order. See Chapter~\ref{ch:content-and-assets}.
\item[.xnb] Real XNA's compiled binary content-pipeline output format. CNA has a real, growing
  binary \texttt{.xnb} reader --- covering \cnaclass{Texture2D}, \cnaclass{SpriteFont}, the
  five stock effects, \cnaclass{SoundEffect}, \cnaclass{Song}, and \cnaclass{Model} at the
  time of writing --- checked first in \cnaclass{ContentManager}'s resolution order, ahead of
  both a direct native-format load and \texttt{.cnj}. See Chapter~\ref{ch:content-and-assets}, and this book's migration
  guide, Chapter~\ref{ch:migration-guide}.
\item[MojoShader] The bytecode-to-GLSL cross-compiler FNA uses to run compiled Direct3D~9
  shader bytecode. CNA has no equivalent, which is the root cause of the shader/.fx bytecode
  gap covered in Chapter~\ref{ch:shader-fx-gap}.
\item[MemoryStream] sharp-runtime's in-memory \cnaclass{IO::Stream} implementation, used where
  a real file path is unavailable (for example content loaded from Android assets). Closing it
  makes I/O operations throw \cnaclass{ObjectDisposedException}, but deliberately preserves the
  bytes for post-close \texttt{ToArray()} / \texttt{GetBuffer()} extraction. Its raw-buffer
  constructor is a known divergence: it creates a read-only stream where real .NET defaults to
  writable. See Chapter~\ref{ch:sharp-runtime-namespaces}.
\item[CNAEXT] CNA's marker convention (a macro, and a matching compile-time build option) for
  any declaration inside the XNA namespace tree that is not part of the real XNA~4.0 API. See
  Chapter~\ref{ch:design-philosophy}.
\item[Oracle (D3D9)] The pixel-for-pixel verification methodology comparing CNA's
  \texttt{DIRECTX9} renderer against a genuinely running XNA~4.0 runtime under Wine, at zero tolerance. See
  Chapter~\ref{ch:direct3d-backends}, and this book's Windows/Wine chapter, Chapter~\ref{ch:windows-wine}.
\item[PbrEffect] A \texttt{CNAEXT} glTF~2.0-style metallic-roughness PBR material effect ---
  unlike the five stock effects, it has no real-XNA counterpart at all, since real XNA~4.0
  predates the glTF PBR model. Available on the WebGPU renderer. See Chapter~\ref{ch:webgpu-backend}.
\item[ReflectiveReader<T>] Real .NET's reflection-based automatic content-reader generation
  mechanism. CNA has no equivalent by permanent, explicit scope decision: only manually
  registered, explicit C\texttt{++} \cnaclass{ContentTypeReader}s are supported, so content
  originally built against an implicit \texttt{ReflectiveReader<T>} needs an explicit
  \texttt{ContentTypeWriter}/reader pair rebuilt and registered before it will load. See
  Chapter~\ref{ch:content-and-assets}.
\item[SDL3] The cross-platform windowing/input/audio library CNA is built on. Not itself a
  graphics renderer --- \texttt{SDL\_RENDERER} is built on
  SDL3's own 2D renderer specifically.
\item[SDL\_GPU] A CNA graphics renderer built directly on SDL3's own \texttt{SDL\_GPU} API, using a
  fixed HLSL-style resource-binding order and a runtime GLSL-to-SPIR-V compile via
  \texttt{libshaderc}. See Chapter~\ref{ch:sdlgpu-backend}.
\item[SDL\_RENDERER] CNA's default renderer outside Linux and Emscripten, selected by
  \texttt{cmake/RendererSelection.cmake};
  a thin, 2D-only wrapper over SDL3's own 2D renderer. See Chapter~\ref{ch:sdl-renderer}.
\item[SharpRuntime / sharp-runtime] The sibling C\texttt{++} reimplementation of a subset of
  the .NET base class library that CNA depends on for its non-XNA-specific vocabulary
  (\texttt{System::Object}, \texttt{TimeSpan}, \texttt{EventHandler<T>}, and more). Covered in
  full in Part~VIII.
\item[SkinnedModelEXT] The CNA-original, \texttt{CNAEXT} real-rendering extension for the
  Avatar system, structurally separate from the faithful, inert base Avatar API. See
  Chapter~\ref{ch:avatar}.
\item[SOFTWARE] A real CPU triangle rasterizer (edge
  functions, perspective-correct interpolation, per-pixel depth test) that never touches SDL
  video/X11/a GPU at all, making it a renderer that can produce a genuine screenshot in a
  headless, display-less environment. See
  Chapter~\ref{ch:graphicsdevice}.
\item[SpriteBatch] The primary 2D rendering API in XNA and CNA alike. See Chapter~\ref{ch:spritebatch}.
\item[Stock effects] The five built-in XNA shader effects --- \cnaclass{BasicEffect},
  \cnaclass{AlphaTestEffect}, \cnaclass{DualTextureEffect}, \cnaclass{EnvironmentMapEffect},
  \cnaclass{SkinnedEffect}. See Chapter~\ref{ch:stock-effects}.
\item[Storage] The \cnans{Microsoft::Xna::Framework::Storage} namespace --- a real, if
  small, XNA namespace for user-data persistence (\cnaclass{StorageDevice},
  \cnaclass{StorageContainer}), backed by a straightforward local-filesystem implementation
  underneath a faithfully-preserved fake-async API shape. See Chapter~\ref{ch:storage}.
\item[VertexDeclaration] The renderer-independent description of a vertex layout: a stride plus
  \cnaclass{VertexElement} records naming offsets, formats, and usages. It is pure data and can
  be constructed on CNA's 2D-only SDL renderer; an actual 3D draw that uses it is where
  that renderer correctly throws. See Chapters~\ref{ch:math-conventions},
  \ref{ch:stock-effects}, and~\ref{ch:shader-fx-gap}.
\item[VULKAN] A CNA graphics renderer targeting the Vulkan API directly. See
  Chapter~\ref{ch:vulkan-backend}.
\item[WEBGPU] CNA's experimental graphics renderer targeting native \texttt{wgpu-native}.
  See Chapter~\ref{ch:webgpu-backend}.
\item[XACT] Microsoft's Cross-platform Audio Creation Tool format/engine, the format CNA's own
  Audio namespace parses and plays back through SDL3\_mixer rather than a FAudio-equivalent
  engine. See Chapter~\ref{ch:audio-system}.
\item[xna4-spec] The sibling repository providing a complete, machine-readable XML
  transcription of the real XNA~4.0 API documentation, used to audit CNA's coverage against a
  ground truth rather than against FNA's own source. See Chapter~\ref{ch:xna4-spec-auditing}.
\end{description}
